% TOP_PROBLEM: {"id":30007548,"problem_number":"LOCAL-30007548","title":"Financing innovation","category_id":21,"category":{"id":21,"name":"economics","display_name":"Economics","description":"Open economic theory statements, published research agendas, and proposed scoped specifications; question provenance and historical priority are qualified per record.","slug":"economics","order_index":21,"created_at":"2026-10-08T00:00:00Z"},"status":"research_question","external_url":null,"legacy_ids":[],"published":false,"statement":"Design incentive-compatible financing and public support for young innovative firms with hidden quality, effort, and innovation spillovers.","economics":{"original_id":37,"field":"Growth and productivity","kind":"design","formalization_basis":"adapted_research_question","source_status":"explicit_open","priority_status":"earliest_found","priority_note":"This is the earliest explicit relevant statement verified in this audit, not a claim of original priority; older literature and earlier versions may contain antecedents.","earliest_verified_reference":{"authors":["Ufuk Akcigit"],"title":"Creative destruction in economic growth","year":2026,"url":"https://onlinelibrary.wiley.com/doi/10.1111/sjoe.70037","doi":"10.1111/sjoe.70037","locator":"Section 7, “Open questions,” item (1)","evidence_summary":"The author explicitly lists financing constraints facing young innovative firms as an open question, stressing limited treatment of financial frictions in growth frameworks."},"current_reference":{"authors":["Ufuk Akcigit"],"title":"Creative destruction in economic growth","year":2026,"url":"https://onlinelibrary.wiley.com/doi/10.1111/sjoe.70037","doi":"10.1111/sjoe.70037","locator":"Section 7, “Open questions,” item (1)","evidence_summary":"The author explicitly lists financing constraints facing young innovative firms as an open question, stressing limited treatment of financial frictions in growth frameworks."},"formalization_note":"Section 7 explicitly identifies young-innovative-firm finance as open. The private-information mechanism and welfare objective are representative adaptations."},"statusQualification":"Published question or research agenda verified at the cited locator. The mathematical specification is adapted for this catalogue and includes additional assumptions; source publication does not establish first-ever priority or certify this exact model remains unsolved. Section 7 explicitly identifies young-innovative-firm finance as open. The private-information mechanism and welfare objective are representative adaptations.","statusReviewedAt":"2026-10-08"}
% FIRST_POSTED: 2026-10-08
% ENTRY_KIND: statement_only
% !TeX program = lualatex
\documentclass[11pt]{article}
\usepackage{fontspec}
\usepackage[a4paper,margin=25mm]{geometry}
\usepackage{amsmath,amssymb,mathtools}
\usepackage{xurl}
\usepackage[hidelinks]{hyperref}
\newcommand{\sref}[2]{\hyperlink{source:#1:#2}{[S#2]}}
\setlength{\emergencystretch}{3em}
\begin{document}
% problemId:30007548
\section*{Financing innovation}\label{30007548}
\subsection{Definitions and mathematical statement}
Consider young firms with privately known research quality \(q\), wealth \(w\), and observable signals \(s\). A project requires funding \(k\), produces an uncertain innovation with specified probability \(p(q,k,e)\), and depends on unobserved effort \(e\). Private cash flows differ from social innovation value because of knowledge spillovers and business stealing. Investors choose contracts contingent on observable outcomes; public policy can provide a bounded guarantee or co-investment budget. Let \(\mathcal C\) be a stated feasible contract family.
Find financing and public-support rules maximizing expected social innovation surplus net of research resources, screening costs, and the financing cost of public funds. Require lender participation, entrepreneur participation, truthful quality revelation where reports are used, and effort incentive compatibility. For each type, a financed project must satisfy the declared equilibrium contract constraints rather than an assumed social-value threshold.
Determine when guarantees or equity co-investment expand valuable radical innovation and when they mainly fund weak projects or displace private finance. Identify the quality distribution, effort response, and spillover parameters using funding discontinuities, contract variation, and independently measured innovation outcomes. Compare policies under common fiscal budgets and report robustness to private-information misspecification. The target is a scoped financing mechanism for young-firm innovation; it does not claim that financial constraints or entrepreneurship have no established theoretical treatment.

\par\textbf{Formulation and scope.} Section 7 explicitly identifies young-innovative-firm finance as open. The private-information mechanism and welfare objective are representative adaptations.

\par\textbf{Open-question provenance.} An explicit question or research agenda was verified at the source locator. This does not by itself certify that every later version remains unresolved \sref{30007548}{1}.

\par\textbf{Historical priority.} This is the earliest explicit relevant statement verified in this audit, not a claim of original priority; older literature and earlier versions may contain antecedents.

\subsection{Short English statement}
Design incentive-compatible financing and public support for young innovative firms with hidden quality, effort, and innovation spillovers.

\subsection{Sources}
\begin{itemize}\raggedright
\item[\textnormal{[S1]}]\hypertarget{source:30007548:1}{} Ufuk Akcigit. 2026. Creative destruction in economic growth. Section 7, “Open questions,” item (1).
\url{https://onlinelibrary.wiley.com/doi/10.1111/sjoe.70037}.
DOI: 10.1111/sjoe.70037.
{\small\textit{Earliest verified question/agenda reference; historical priority qualified below. The author explicitly lists financing constraints facing young innovative firms as an open question, stressing limited treatment of financial frictions in growth frameworks.}}
\end{itemize}
\end{document}
